\chapter{Chapter08 — Product LTF Noise}
%%% generated by the blueprint pipeline from FABL.Chapter08.ProductLTFNoise
\section{Overview}

Declarations of the Lean module \texttt{FABL.Chapter08.ProductLTFNoise}.

\section{Declarations}

\begin{definition}[Product linear threshold function]
\lean{FABL.IsProductLinearThreshold}
\statementsource{Analysis-of-Boolean-Functions-by-Ryan-ODonnell}{
  pdf-ffbe9c7b3306-p044-b005,
  pdf-ffbe9c7b3306-p113-b003,
  pdf-ffbe9c7b3306-p241-b003
}
\label{FABL.IsProductLinearThreshold}
\leanok
\uses{FABL.Sign, FABL.thresholdSign}
Let $n$ be a natural number, let $\Omega$ be a finite type, and let
$f:\Omega^n\to\{\pm1\}$. The function $f$ is a product linear threshold if there are
real weights $a_{i}(\omega)$ for each coordinate $i$ and symbol $\omega$, and a real
bias $a_0$, such that for every input $x$, $f(x)$ is the threshold sign of
$a_0+\sum_{i=1}^n a_i(x_i)$.
\end{definition}

\begin{definition}[Coordinatewise product input selection]
\lean{FABL.selectProductInput}
\statementsource{Analysis-of-Boolean-Functions-by-Ryan-ODonnell}{
  pdf-ffbe9c7b3306-p034-b004,
  pdf-ffbe9c7b3306-p241-b003
}
\label{FABL.selectProductInput}
\leanok
Let $n$ be a natural number, let $\Omega$ be a type, and let $u,v\in\Omega^n$. For a
sign vector $x\in\{-1,1\}^n$, the selected input is the vector whose $i$th coordinate is
$u_i$ when $x_i=1$ and $v_i$ otherwise.
\end{definition}

\begin{theorem}[Two-point selection identity]
\lean{FABL.selectedScore_eq_midpoint_add_difference}
\label{FABL.selectedScore_eq_midpoint_add_difference}
\leanok
\uses{FABL.Sign, FABL.signValue}
\difficulty{4}
Let $\Omega$ be a type, let $a:\Omega\to\mathbb{R}$, choose
$\text{plus},\text{minus}\in\Omega$, and let $x\in\{−1,1\}$. Then $a(\text{plus})$ is
selected when $x=1$ and $a(\text{minus})$ otherwise, and the selected value is
$\frac{a(\text{plus})+a(\text{minus})}{2}+\frac{a(\text{plus})-a(\text{minus})}{2}x$.
\end{theorem}

\begin{theorem}[Restriction to selected product inputs]
\lean{FABL.isLinearThreshold_selectProductInput}
\label{FABL.isLinearThreshold_selectProductInput}
\leanok
\uses{FABL.IsLinearThreshold, FABL.IsProductLinearThreshold, FABL.Sign, FABL.selectProductInput, FABL.selectedScore_eq_midpoint_add_difference}
\difficulty{4}
Let $\Omega$ be a finite set, let $f:\Omega^n\to\{-1,1\}$ be a linear threshold
function, and choose two elements of $\Omega$ at each coordinate, denoted $u_i^+$ and
$u_i^-$. The function on $\{-1,1\}^n$ that evaluates $f$ at the input whose $i$th
coordinate is $u_i^+$ when $x_i=1$ and $u_i^-$ otherwise is a linear threshold function.
\end{theorem}

\begin{definition}[Boolean sample selected by a sign]
\lean{FABL.boolSampleSelectedBySign}
\statementsource{Analysis-of-Boolean-Functions-by-Ryan-ODonnell}{pdf-ffbe9c7b3306-p241-b003}
\label{FABL.boolSampleSelectedBySign}
\leanok
\uses{FABL.Sign}
For Boolean values $u,v$ and a sign $s\in\{\pm1\}$, the selected Boolean value is $u$
when $s=1$ and $v$ otherwise.
\end{definition}

\begin{definition}[Selector agreement at zero]
\lean{FABL.printedSelectorAgreementAtZero}
\label{FABL.printedSelectorAgreementAtZero}
\leanok
\uses{FABL.Sign, FABL.boolSampleSelectedBySign, FABL.pmfExpectation, FABL.uniformPMF}
Choose independent uniform bits $a,b\in\{0,1\}$ and independent uniform signs
$x,y\in\{\pm1\}$. The printed selector agreement at zero is the probability that the
Boolean samples selected from $a,b$ by $x$ and $y$ agree.
\end{definition}

\begin{definition}[Product Boolean agreement at zero]
\lean{FABL.productBoolAgreementAtZero}
\label{FABL.productBoolAgreementAtZero}
\leanok
\uses{FABL.pmfExpectation, FABL.productCoordinateNoisePMF, FABL.uniformPMF}
Choose a uniform bit $x\in\{0,1\}$, then draw $y$ from the product-coordinate noise
distribution at parameter $0$ conditional on $x$. The product Boolean agreement at zero
is the probability that $x=y$.
\end{definition}

\begin{theorem}[Selector agreement at zero]
\lean{FABL.printedSelectorAgreementAtZero_eq_three_fourths}
\label{FABL.printedSelectorAgreementAtZero_eq_three_fourths}
\leanok
\uses{FABL.boolSampleSelectedBySign, FABL.expect_sign, FABL.pmfExpectation_uniformPMF_eq_expect, FABL.printedSelectorAgreementAtZero}
\difficulty{3}
Choose two independent uniform bits $U,V$ and two independent uniform signs
$X,Y\in\{-1,1\}$. For each sign, select $U$ if it is $1$ and $V$ otherwise. The
probability that the selections made using $X$ and $Y$ agree, averaged over all four
choices, is $3/4$.
\end{theorem}

\begin{theorem}[Agreement under zero noise]
\lean{FABL.productBoolAgreementAtZero_eq_one_half}
\label{FABL.productBoolAgreementAtZero_eq_one_half}
\leanok
\uses{FABL.pmfExpectation_productCoordinateNoisePMF, FABL.pmfExpectation_uniformPMF_eq_expect, FABL.productBoolAgreementAtZero}
\difficulty{3}
For a uniformly distributed Boolean variable, choose a second Boolean variable by the
coordinate noise rule with retention probability zero. The probability that the two
variables agree is $1/2$.
\end{theorem}

\begin{theorem}[Selector and product coupling disagree]
\lean{FABL.printed_product_selector_coupling_false}
\label{FABL.printed_product_selector_coupling_false}
\leanok
\uses{FABL.printedSelectorAgreementAtZero, FABL.printedSelectorAgreementAtZero_eq_three_fourths, FABL.productBoolAgreementAtZero, FABL.productBoolAgreementAtZero_eq_one_half}
\difficulty{3}
Choose independent uniform bits $U,V$ and independent uniform signs $X,Y\in\{-1,1\}$.
Select $U$ for sign $+1$ and $V$ for sign $-1$, and let $A$ be the probability that the
selections using $X$ and $Y$ agree. Separately, choose a uniform bit and independently
resample a uniform bit; let $B$ be the probability that they agree. Then $A\ne B$.
\end{theorem}
